\section{Lecture 23 (November 26th)}
\begin{thm}
We continue to calculate the decay rate from a $2p$ state to a $1s$ state emiting a photon. We use Fermi's golden rule, where
\[\mathrm{\Gamma} =\dfrac{2\pi }{\hbar }\sum_{\mathrm{final}} |\langle f|M|i\rangle |^2\delta (\mathrm{\Sigma} E)\]
We've expressed all states as a tensor product of the atomic wavefunction and the photon wavefunction, and 2nd quantised $\vec{A}$. The transition Hamiltonian is given as
\[M=\dfrac{e}{m_{e}}\vec{A}\cdot \vec{p}\]
that came from the first order of the original Hamiltonian. To compute the inner product of the initial and final atomic states, we needed a particular trick, where we needed to calculate
\[\begin{align*}
\langle \phi _{f}|\vec{x}\cdot \hat{\epsilon }|\phi _{i}\rangle =&\int d\mathrm{\Omega} \,Y^{*}_{l_{f}m_{f}}(\hat{\epsilon }\cdot \hat{r})Y_{l_{i}m_{i}}=\int d\mathrm{\Omega} \,Y^{*}_{l_{f}m_{f}}(\epsilon_{x}r_{x}+\epsilon _{y}r_{y}+\epsilon _{z}r_{z})Y_{l_{i}m_{i}}
\end{align*}\]
where
\[\epsilon_{x}r_{x}+\epsilon_{y}r_{y}+\epsilon _{z}r_{z}=\sqrt{\dfrac{4\pi }{3}}\Big(\epsilon_{z}Y_{10}+\dfrac{-\epsilon_{x}+i\epsilon_{y}}{\sqrt{2}}Y_{11}+\dfrac{\epsilon _{x}+i\epsilon _{y}}{\sqrt{2}}Y_{1,-1}\Big)\]
noting that
\[\begin{align*}
(Y_{lm})^{*}=(-1)^{m}Y_{l,-m}\qquad\&\qquad \mathcal{P}Y_{lm}\rightarrow (-1)^{l}Y_{lm}
\end{align*}\]
where $\mathcal{P}$ is the parity operator. For simplicity, let's assume $\vec{R}\,||\,\hat{z}$. $\epsilon _{z}=0$ and $m_{f}-m_{i}=\pm 1$. We can then compute the above using
\[\int d\mathrm{\Omega} \,Y^{*}_{l'm'}Y_{lm}=\delta _{ll'}\delta _{mm'}\]
Through all this, we get to obtain
\[d\mathrm{\Gamma} =\dfrac{1}{3}\sum _{m_{i}=\pm 1,0}d\mathrm{\Gamma} _{m} \propto (\epsilon_{x}^2+\epsilon _{y}^2+\epsilon_{z}^2)\]
and numerically, $\mathrm{\Gamma} (2p\rightarrow 1s)=0.62\times 10^{6}$ per second. 
\end{thm}
\vspace{2ex}
\begin{ex}
To this point, we have considered no change in spin, $\mathrm{\Delta }S=0$. However, $\mathrm{\Delta }S\ne 0$ is indeed possible, where 
\[\dfrac{eg}{m_{e}c}\vec{S}\cdot \vec{B}\]
and $\vec{B}=\nabla \times \vec{A}$. Consider the reaction:
\[\gamma +d\rightarrow p+n\]
where $d$ is in the ${}^{3}S_{1}$ triplet state and $p+n$ is in the ${}^{1}S_{0}$ singlet state. Thus, $\mathrm{\Delta }S=1$. We find
\[\langle f|(g_{p}\vec{S}_{p}+g_{n}\vec{S}_{n})\cdot \vec{B}|i\rangle \ne 0\]
in the homework problems.
\end{ex}
\vspace{2ex}
\begin{thm}
(No clone theorem) Consider an arbitrary qubit $|x\rangle \in {\bm C}^2$. Consider a clone map (a copy map):
\[G:{\bm C}^2\times {\bm C}^2\rightarrow {\bm C}^2\times {\bm C}^2\]
which is a unitary operator. It would send
\[|x\rangle \otimes |0\rangle \rightarrow |x\rangle \otimes |x\rangle \]
where we emphasize that we do not know $|x\rangle $. The no clone theorem suggests that there doesn't exist such a $G$ for a general $|x\rangle $. 
\end{thm}
\vspace{2ex}
 
